\documentclass{article}
\usepackage[utf8]{inputenc}
\usepackage{amsmath}
\usepackage{geometry}
\geometry{margin=2cm}
\begin{document}

\section*{Análise da Questão 137 — ENEM 2024 — Matemática}

\subsection*{Enunciado}
Ao calcular a média de suas notas em 4 provas, um estudante dividiu, por engano, a soma das notas por 5. Com isso, a média obtida foi 1 unidade menor do que deveria ser, caso fosse calculada corretamente. O valor correto da média das notas desse estudante é:

\subsection*{Solução Técnica}

Para determinar a média correta das 4 notas do estudante, siga os passos:

\begin{enumerate}
  \item Defina a soma das notas como $S$.
  \item A média correta é
  \begin{equation}
  M = \frac{S}{4}.
  \end{equation}
  \item O estudante calculou a média dividindo a soma por 5, obtendo a média errada:
  \begin{equation}
  M_{errada} = \frac{S}{5}.
  \end{equation}
  \item Sabe-se que a média errada foi 1 unidade menor:
  \begin{equation}
  M - M_{errada} = 1.
  \end{equation}
  \item Substituindo $M$ e $M_{errada}$:
  \begin{equation}
  \frac{S}{4} - \frac{S}{5} = 1.
  \end{equation}
  \item Calculando o mínimo múltiplo comum $20$ e resolvendo:
  \begin{equation}
  \frac{5S - 4S}{20} = 1 \implies \frac{S}{20} = 1 \implies S = 20.
  \end{equation}
  \item Calculando a média correta:
  \begin{equation}
  M = \frac{20}{4} = 5.
  \end{equation}
\end{enumerate}

\textbf{Resposta final:} $\boxed{5}$.

\subsection*{Explicação Didática}

O estudante cometeu o erro de dividir a soma das notas por 5 em vez de 4, fazendo com que a média obtida fosse 1 unidade menor do que a correta. Montando uma equação com essa condição, descobrimos que a média certa é 5.

\subsection*{Passo a Passo}

\begin{enumerate}
  \item Defina a soma das 4 notas como $S$ e a média correta como $M = \frac{S}{4}$.
  \item A média calculada de forma errada foi $\frac{S}{5}$.
  \item Sabemos que a diferença entre a média correta e a errada é 1:
  \[
    M - \frac{S}{5} = 1.
  \]
  \item Substituindo $M$ por $\frac{S}{4}$ temos:
  \[
    \frac{S}{4} - \frac{S}{5} = 1.
  \]
  \item Encontramos o mínimo múltiplo comum para as frações e resolvemos:
  \[
    \frac{5S - 4S}{20} = 1 \Rightarrow \frac{S}{20} = 1 \Rightarrow S = 20.
  \]
  \item A média correta é então
  \[
    M = \frac{S}{4} = \frac{20}{4} = 5.
  \]
\end{enumerate}

\subsection*{Dificuldades Comuns}
\begin{itemize}
  \item Confundir o número correto de provas para dividir ao calcular a média.
  \item Não montar uma equação para relacionar as médias correta e incorreta.
  \item Erros no cálculo com frações e manipulação algébrica.
  \item Confundir soma total das notas com o valor da média.
\end{itemize}

\subsection*{Dicas para Ensino}
\begin{itemize}
  \item Reforce o conceito de média como soma das notas dividida pelo número de provas.
  \item Incentive a criação de equações a partir das informações dadas.
  \item Pratique manipulação de frações e equações lineares simples.
  \item Use exemplos com erros comuns para reforçar a compreensão.
\end{itemize}

\subsection*{Alternativas Incorretas}

\begin{itemize}
  \item \textbf{4}: Confusão entre denominador e média, parecendo próximo da média mas não atende o enunciado.
  \item \textbf{6}: Supõe adicionar 1 na média errada sem respeitar a relação correta.
  \item \textbf{19}: Erro no cálculo e interpretação, número muito alto para média.
  \item \textbf{21}: Confunde soma total com média, não é o valor pedido.
\end{itemize}

\subsection*{Perfil da Questão}

\begin{itemize}
  \item Habilidade principal: cálculo de média e manipulação algébrica.
  \item Habilidades secundárias: interpretação e resolução de equações simples.
  \item Exigência cognitiva: média.
  \item Dificuldade: fácil.
  \item Requer modelagem algébrica e cálculo.
\end{itemize}

\subsection*{Revisão de Consistência}
O enunciado, solução e explicações estão coerentes e consistentes. A resposta final é correta, com raciocínio lógico adequado e abordagens dos erros comuns nas alternativas e explicações. Não foram identificadas inconsistências importantes.


\end{document}